%% RESUMO EM LINGUA ESTRANGEIRA (elemento obrigatorio -- NBR 6028)
%% As keywords sao definidas em config/dados.tex.

\begin{abstract}
Large language models are being used in public administration to classify
documents, yet they do not eliminate errors, and their adoption requires
distinguishing the cases in which automated classification is reliable from
those in which human review remains necessary. This dissertation investigates
to what extent uncertainty signals produced by the models themselves make it
possible to decide when a classification can be accepted automatically and
when it should be referred to a person. The theoretical framework draws on
selective classification, which trades coverage for accuracy, and on the
literature on self-evaluation and calibration of language models. The
research, quantitative and applied, compares the classification proposed by
the models with an official classification assigned by experts to a
public-sector document base and evaluates each signal through risk-coverage
curves. It aims to provide a verifiable criterion for referral to human
review, consistent with the requirement of oversight proportional to risk.
\end{abstract}
